% ==============================================================================
% SECCION 4: DISCUSION Y ANALISIS FISICO DE ERROR
% ==============================================================================
\section{Discusión y Análisis Físico de Error}\label{sec:discusion}

\subsection{Comparativa Cuantitativa: Teórico/CAD vs. Simulación vs. Hardware}\label{subsec:comparacion_general}

La Tabla~\ref{tab:comparacion_general} confronta los parámetros físicos proyectados desde el modelo CAD (FreeCAD~1.1) y hojas de datos frente a la simulación no lineal a \qty{1.0}{\kilo\hertz} y los registros experimentales en hardware real.

\begin{table}[!t]
\centering
\caption{Comparativa Cuantitativa: Modelo Teórico/CAD, Simulación (\qty{1}{\kilo\hertz}) y Hardware Real (\texttt{COM13}).}
\label{tab:comparacion_general}
\resizebox{\columnwidth}{!}{%
\begin{tabular}{@{}lcccc@{}}
\toprule
\textbf{Parámetro / Métrica} & \textbf{Teór./CAD} & \textbf{Sim.} & \textbf{Real} & \textbf{Err.\ Rel.} \\
\midrule
$K_v$ [\si{(\radian\per\second)\per\volt}] & $3.40$ & $3.372$ & $3.372$ & $-0.8\%$ \\
$\tau_m$ [\si{\second}] & $0.1200$ & $0.1175$ & $0.1175$ & $-2.1\%$ \\
$\omega_n$ [\si{\radian\per\second}] & $8.95$ & $8.8484$ & $8.8484$ & $-1.1\%$ \\
$\zeta$ [adim.] & $0.0500$ & $0.0649$ & $0.0649$ & $+29.8\%$ \\
$K_{\text{acople}}$ [adim.] & $0.850$ & $0.8737$ & $0.8737$ & $+2.8\%$ \\
$J_{1,\text{piv}}$ [\si{\kilo\gram\square\meter}] & $2.100\times 10^{-4}$ & $2.044\times 10^{-4}$ & $2.044\times 10^{-4}$ & $-2.7\%$ \\
$u_{\text{breakward}}^{-}$ [\si{\volt}]$^*$ & $-1.20$ & $-2.52$ & $-2.52$ & $-110.0\%$ \\
\midrule
$t_{\text{swing}}$ [\si{\second}] & $\le 1.50$ & $1.48$ & $0.98$ & $-33.8\%$ \\
$\mathrm{RMS}(\tilde{\alpha})$ [\si{\degree}] & --- & $1.15$ & $1.99$ & $+73.0\%$ \\
$\mathrm{RMS}(\theta)$ [\si{\degree}] & --- & $4.85$ & $9.31$ & $+92.0\%$ \\
$\mathrm{RMS}(u)$ [\si{\volt}] & --- & $0.58$ & $0.73$ & $+25.9\%$ \\
Caídas en equilibrio & $0$ & $0$ & $0$ & $0.0\%$ \\
\bottomrule
\multicolumn{5}{@{}l@{}}{\scriptsize $^*$Valor nominal simétrico ($\pm 10\%$). Error rel.: bloque sup.\ vs.\ CAD; bloque inf.\ vs.\ Sim.}
\end{tabular}%
}
\end{table}

Los parámetros dimensionales, inerciales y electromecánicos ($K_v$, $\tau_m$, $\omega_n$, $K_{\text{acople}}$ y $J_{1,\text{piv}}$) exhiben una concordancia notable con errores relativos inferiores al \qty{3}{\percent} (máximo $\qty{2.8}{\percent}$), validando la fidelidad geométrica del modelado CAD y la manufactura DfAM sin soportes (Sec.~\ref{subsec:dfam}). Por el contrario, los términos disipativos y no lineales evidencian dos discrepancias severas respecto a las hipótesis ideales: el factor de amortiguamiento $\zeta$ aumenta $+29.8\%$ debido a la fricción seca de los sellos del rodamiento \texttt{608ZZ} y la resistencia aerodinámica del contrapeso, mientras que $u_{\text{breakward}}^{-}$ alcanza $-\qty{2.52}{\volt}$ ($-110.0\%$ respecto al umbral simétrico nominal de $-\qty{1.20}{\volt}$) por asimetría interna de los pares Darlington del puente~H \texttt{L298N}. En lazo cerrado (Fig.~\ref{fig:sim_vs_exp}), el prototipo real completa el \textit{Swing-Up} un \qty{33.8}{\percent} más rápido que la simulación porque la fricción estática inicial suprime un semiciclo de baja amplitud antes del bombeo resonante; en régimen permanente, la holgura de engranajes, el ruido de cuantización y la rigidez del cableado elevan $\mathrm{RMS}(\tilde{\alpha})$, $\mathrm{RMS}(\theta)$ y $\mathrm{RMS}(u)$ sin comprometer la estabilidad ($0$~caídas).

\subsection{Diagnóstico de Falla de las Ganancias Offline (L-BFGS-B)}\label{subsec:diagnostico_lbfgsb}

Las ganancias obtenidas mediante optimización numérica L-BFGS-B fuera de línea (\texttt{analizar\_y\_optimizar\_}\allowbreak\texttt{sistema\_furuta.py}) sobre el modelo simplificado sin holgura, cuantización ni resonancia ($K_{p,\alpha}=52.01$, $K_{i,\alpha}=16.06$, $K_{d,\alpha}=2.23$; $K_{p,\theta}=-2.73$, $K_{i,\theta}=-2.00$, $K_{d,\theta}=-2.58$) provocaron inestabilidad divergente inmediata al transferirse al prototipo STM32 real, excitando una oscilación violenta gobernada por polos en el semiplano derecho en $s_{1,2} = +6.48 \pm 32.91j\text{ rad/s}$ ($\sim\qty{5.24}{\hertz}$, $|z_{1,2}|=1.0065>1$). Tres mecanismos físico-matemáticos explican esta falla:

\begin{enumerate}
    \item \textbf{Competencia integral en sistema SIMO ($K_{i,\alpha}>0$):} Al disponer de un único actuador $u(t)$ para regular dos variables acopladas ($\tilde{\alpha},\theta$), es físicamente imposible anular ambos errores en estado estacionario ante cualquier sesgo de verticalidad ($\sim\qty{0.2}{\degree}$). El equilibrio físico requiere $\tilde{\alpha}_{\text{eq}}\neq 0$, haciendo que $K_{i,\alpha}=16.06$ y $K_{i,\theta}=-2.00$ integren en oposición con deriva sostenida del brazo; simultáneamente, el polo integrador en $s=0$ introduce $-\qty{90}{\degree}$ de retardo de fase en baja frecuencia, erosionando el margen de fase requerido para contener el polo inestable en $s_4 = +\qty{8.293}{\radian\per\second}$ (Sec.~\ref{subsec:euler_lagrange}).
    \item \textbf{Ganancia derivativa excesiva del brazo ($K_{d,\theta}=-2.58$) y $K_{p,\alpha}$ alto:} En la dinámica acoplada $\ddot{\tilde{\alpha}} = \omega_n^2\tilde{\alpha} - 2\zeta\omega_n\dot{\tilde{\alpha}} - K_{\text{acople}}\ddot{\theta}$ con $\ddot{\theta} = (-\dot{\theta}+K_v u)/\tau_m$, una magnitud elevada de $|K_{d,\theta}|=2.58$ frente a $K_{d,\alpha}=2.23$ junto con $K_{p,\alpha}=52.01$ ($3.6\times$ superior al diseño final) amplifican el ruido derivativo de 12~bits y realimentan con signo desestabilizador el retardo de fase introducido por la holgura mecánica, desencadenando la oscilación divergente a \qty{5.24}{\hertz}.
    \item \textbf{Ausencia de atenuación espectral en \qty{11.0}{\hertz}:} Sin filtrado selectivo \textit{Dual-Notch} ni acotamiento de pendiente ($\mathrm{d}u/\mathrm{d}t$), el contenido de alta frecuencia realimentado excita continuamente el ciclo límite propio de la reductora planetaria $80:1$ en \qty{11.0}{\hertz}, saturando el puente~H.
\end{enumerate}

La arquitectura \texttt{M4 Ultra-Silk}~\eqref{eq:m4_dual}--\eqref{eq:dual_notch} resuelve simultáneamente los tres mecanismos: impone $K_{i,\alpha}=0$ estricto para delegar la compensación de sesgos exclusivamente al integrador acotado del brazo ($K_{i,\theta}=-0.75$); reduce $K_{p,\alpha}$ a $14.50$ y modera la acción del brazo a $K_{d,\theta}=-0.85$ bajo saturación subordinada $u_\theta \in [-\qty{2.2}{\volt}, +\qty{2.2}{\volt}]$ para priorizar la verticalidad; e intercala la cascada LPF~IIR ($f_c\approx\qty{15.8}{\hertz}$) junto al filtro \textit{Dual-Notch} (\qty{11.0}{\hertz}/\qty{6.2}{\hertz}) y el \textit{slew rate} adaptativo ($\qty{0.065}{}\to\qty{0.40}{\volt\per\milli\second}$), reubicando los cinco polos en~\eqref{eq:polos_m4} dentro del círculo unitario ($\rho = 0.99934 < 1$).

\subsection{Fuentes Físicas de No Idealidad en Hardware}\label{subsec:fuentes_no_idealidad}

Las discrepancias remanentes entre el modelo idealizado y el banco físico derivan de seis fenómenos estructurales:
\begin{enumerate}
    \item \textbf{Holgura (\textit{backlash}) de la reductora planetaria 80:1:} El juego angular entre dientes del tren epicicloidal introduce un retardo no lineal en cada inversión de par, excitando un ciclo límite en \qty{11.0}{\hertz} suprimido por el Notch~1 en \eqref{eq:dual_notch}.
    \item \textbf{Asimetría BJT Darlington del L298N:} Las caídas internas de saturación ($V_{\text{sat}}\approx\qty{2.0}{\volt}$) y el desbalance entre transistores NPN/PNP provocan umbrales de arranque asimétricos de $+10.0\%$ ($+\qty{1.20}{\volt}$) y $-21.0\%$ ($-\qty{2.52}{\volt}$), compensados suavemente mediante \eqref{eq:comp_l298n}.
    \item \textbf{Ruido de cuantización por diferenciación a \qty{1.0}{\kilo\hertz}:} Con resolución distal de 12~bits ($\Delta\alpha_{\min}=\qty{1.534e-3}{\radian}$), un salto de $1\,\mathrm{LSB}$ en $T_s=\qty{1.0}{\milli\second}$ induce picos derivativos de:
    \begin{equation}
        \Delta\dot{\alpha}_{\text{ruido}} = \frac{1.534\times 10^{-3}}{10^{-3}} = \pm\qty{1.534}{\radian\per\second},
        \label{eq:ruido_derivativo}
    \end{equation}
    atenuados a $-40\text{ dB/década}$ por la cascada LPF ($f_c\approx\qty{15.8}{\hertz}$), mientras el sensor base montado en el rotor previo al reductor multiplica $80\times$ la resolución de $\theta$ ($\qty{0.0011}{\degree}$).
    \item \textbf{Fricción mixta viscosa + Coulomb del rodamiento 608ZZ:} La fricción seca estática de sus sellos ($c_c=\qty{0.6179}{\radian\per\second\squared}$ en \eqref{eq:friccion_mixta}) eleva $\zeta$ en $+29.8\%$ e impide un cero matemático absoluto, fijando un residuo de $\mathrm{RMS}(\tilde{\alpha})=\qty{1.99}{\degree}$.
    \item \textbf{Par elástico residual del cableado distal:} Los cuatro conductores del bus $\text{I}^2\text{C}_2$ del AS5600 distal ejercen un par restaurador dependiente del ángulo, compensado como sesgo de posición $\theta$ por el integrador de $\text{PID}_\theta$.
    \item \textbf{Desacoplamiento por inercia reflejada:} La razón constructiva $J_{\text{ref}}/(m_1 L_0 l_1) \approx 103.4 \gg 1$ (Sec.~\ref{subsec:euler_lagrange}) atenúa por debajo del \qty{1}{\percent} la retroacción del péndulo sobre el brazo, desacoplando eficazmente la planta motriz ($R^2=98.51\%$).
\end{enumerate}
